% Source: https://artofproblemsolving.com/community/c3h2668799_identifying_sides_for_trig
% Stable identity: aops-topic-2668799-q01
\textbf{(a)} Relative to the acute angle $\theta$: the \emph{hypotenuse} is the side opposite the right angle; the \emph{opposite} leg is the leg not touching $\theta$; the \emph{adjacent} leg is the leg that forms $\theta$ together with the hypotenuse. Then
\[\sin\theta=\dfrac{\text{opposite}}{\text{hypotenuse}},\qquad\cos\theta=\dfrac{\text{adjacent}}{\text{hypotenuse}},\qquad\tan\theta=\dfrac{\text{opposite}}{\text{adjacent}}.\]
The student's ratio is $\tan\theta$, not $\sin\theta$.

\textbf{(b)} The hypotenuse is $\sqrt{3^2+4^2}=5$, so
\[\sin\theta=\boxed{\dfrac35},\qquad\cos\theta=\boxed{\dfrac45},\qquad\tan\theta=\boxed{\dfrac34}.\]

\begin{asy}
usepackage("amsmath");
size(4.5cm);
defaultpen(fontsize(10pt));
pair A=(0,0),B=(4,0),C=(4,3);
draw(A--B--C--cycle,linewidth(0.8));draw(B+(-0.3,0)--B+(-0.3,0.3)--B+(0,0.3));
label("$4$ (adjacent)",(A+B)/2,S,fontsize(8pt));label("$3$ (opposite)",(B+C)/2,E,fontsize(8pt));label("$5$ (hypotenuse)",(A+C)/2,NW,fontsize(8pt));
label("$\theta$",A,dir(18)*4);
\end{asy}
